\chapter{Sicurezza dei microservizi}
\label{cap:sicurezza-micro}

\begin{inbreve}
Dividere un'applicazione in microservizi cambia il modo di proteggerla: più punti d'ingresso, più comunicazioni in
rete, più team con pratiche diverse. Il capitolo elenca prima le sette sfide della sicurezza nei microservizi. Poi
presenta, con lo stesso schema del capitolo~\ref{cap:smell}, gli \emph{smell} che indicano possibili violazioni della
sicurezza e i refactoring che li mitigano. Si chiude con il procedimento per eliminarli. Presuppone i
capitoli~\ref{cap:microservizi}, \ref{cap:smell} e~\ref{cap:sicurezza}.
\end{inbreve}

\section{Le sfide}
\label{sec:secm-sfide}

\subsection{Più ampia è la superficie d'attacco, più alto è il rischio}

La \termine{superficie d'attacco} è l'insieme dei punti da cui un attaccante può tentare di entrare nel sistema.
\begin{itemize}
  \item In un \textbf{monolite} i componenti comunicano all'interno di un unico processo.
  \item Nei \textbf{microservizi} i servizi comunicano con chiamate remote e i punti d'ingresso possono essere molti.
    La superficie d'attacco si allarga.
\end{itemize}
\uline{La sicurezza di un'applicazione è pari alla robustezza del suo punto più debole.}

\subsection{I controlli distribuiti pesano sulle prestazioni}

\begin{itemize}
  \item In un \textbf{monolite} il controllo di sicurezza (\ingl{security screening}) si fa una volta sola; poi la
    richiesta è inoltrata al componente giusto.
  \item Nei \textbf{microservizi} ogni servizio deve fare un controllo di sicurezza indipendente. Può dover contattare
    un \emph{security token service} remoto, il servizio che emette e verifica i token di sicurezza. Controlli ripetuti
    e distribuiti pesano sulle prestazioni.
\end{itemize}
Le soluzioni sono due.
\begin{itemize}
  \item Un ripiego: \emph{trust-the-network}. Ci si fida della rete interna e si saltano i controlli in ogni
    microservizio.
  \item La tendenza dell'industria: i principi dello \termine{zero-trust networking}, in cui nessuna richiesta è
    considerata fidata solo perché arriva dalla rete interna. In questo caso bisogna tenere conto delle prestazioni
    complessive.
\end{itemize}

\subsection{Stabilire la fiducia tra microservizi richiede automazione}

La comunicazione tra servizi deve avvenire su canali protetti. Se si usano i certificati (\rif{sec:sic-tls}):
\begin{itemize}
  \item ogni microservizio deve ricevere un certificato, con la relativa chiave privata, con cui autenticarsi presso
    gli altri microservizi;
  \item il microservizio che riceve una chiamata deve sapere come validare il certificato del chiamante;
  \item serve un modo per stabilire all'inizio (\ingl{bootstrap}) la fiducia tra microservizi;
  \item bisogna anche revocare i certificati e sostituirli periodicamente (\ingl{rotate}).
\end{itemize}
Per gestire distribuzioni su larga scala, con centinaia di microservizi, serve l'\textbf{automazione}, per esempio con
Docker e Kubernetes (capitolo~\ref{cap:kubernetes}).

\subsection{Tracciare richieste che attraversano più microservizi è difficile}

\begin{itemize}
  \item Un \termine{log} è la registrazione di un evento di un servizio.
  \item I log si possono aggregare per produrre metriche che riflettono lo stato del sistema, per esempio il numero
    medio di richieste di accesso non valide all'ora. Le metriche possono far scattare allarmi.
  \item Le \termine{trace} seguono una richiesta dal punto in cui entra nel sistema al punto in cui ne esce.
\end{itemize}
A differenza di quanto accade nei monoliti, una richiesta può attraversare più microservizi. Mettere in relazione le
richieste tra microservizi diversi è quindi difficile. Ci sono strumenti appositi: per esempio Prometheus e Grafana
per monitorare le richieste in arrivo, Jaeger e Zipkin per il \emph{distributed tracing}.

\subsection{I container complicano credenziali e politiche}

I container sono \emph{server immutabili}: dopo l'avvio non cambiano stato (\rif{sec:k8s-immutabile}). Questo
semplifica molto la distribuzione e la scalabilità orizzontale. Però:
\begin{itemize}
  \item per ogni servizio bisogna mantenere una lista \emph{dinamica} dei client ammessi e un insieme \emph{dinamico}
    di politiche di controllo degli accessi. Per esempio si scaricano le politiche aggiornate da un endpoint di
    amministrazione delle politiche, con un modello push o pull;
  \item ogni servizio deve mantenere anche le proprie credenziali, da sostituire periodicamente. Per esempio si tengono
    le credenziali nel filesystem del container e si iniettano all'avvio.
\end{itemize}

\subsection{La distribuzione rende difficile condividere il contesto utente}

Il \emph{contesto utente} (chi è l'utente e quali attributi ha) va passato in modo esplicito da un microservizio
all'altro. La sfida è costruire una fiducia tra microservizi tale che chi riceve accetti il contesto utente passato dal
chiamante.

Una soluzione diffusa è il \termine{JSON Web Token} (JWT): messaggi che trasportano gli attributi dell'utente in modo
crittograficamente sicuro. I messaggi sono firmati, così chi li riceve può verificare che non sono stati alterati.

\subsection{Le responsabilità di sicurezza sono divise tra team}

\begin{itemize}
  \item In un'architettura poliglotta squad diverse (\rif{sec:devops-spotify}) possono usare stack tecnologici diversi.
  \item Team diversi possono seguire buone pratiche di sicurezza diverse e usare strumenti diversi per il security
    testing statico e dinamico.
  \item Le responsabilità di sicurezza sono quindi distribuite tra più team.
\end{itemize}
Le organizzazioni adottano spesso un approccio ibrido: un team centrale per la sicurezza ed esperti di sicurezza
dentro i team.

\begin{idea}[Le sette sfide]
\begin{enumerate}
  \item Più ampia è la superficie d'attacco, più alto è il rischio.
  \item I controlli distribuiti pesano sulle prestazioni.
  \item Stabilire la fiducia tra microservizi richiede automazione.
  \item Tracciare richieste che attraversano più microservizi è difficile.
  \item I container complicano la gestione di credenziali e politiche.
  \item La distribuzione rende difficile condividere il contesto utente.
  \item Le responsabilità di sicurezza sono divise tra team diversi.
\end{enumerate}
\end{idea}

Riferimento: P. Siriwardena, N. Dias, \emph{Microservices Security in Action}, 2020.

\section{Smell e refactoring per la sicurezza}
\label{sec:secm-smell}

Come per i principi di design (\rif{sec:smell-problema}), una \emph{multivocal review} di letteratura bianca e grigia,
con 58 riferimenti, risponde a due domande.
\begin{enumerate}
  \item Quali sono gli smell più riconosciuti che indicano possibili violazioni della sicurezza nelle applicazioni a
    microservizi?
  \item Come si rifattorizza un'applicazione a microservizi per mitigare l'effetto di uno smell di sicurezza?
\end{enumerate}
Il risultato è una tassonomia con \textbf{10 refactoring} associati a \textbf{10 smell}, che possono colpire
\textbf{3 proprietà di sicurezza}.

\begin{definizione}[Proprietà di sicurezza]
Le tre proprietà vengono dallo standard di qualità del software ISO/IEC 25010.
\begin{description}[leftmargin=0pt, style=unboxed]
  \item[Confidentiality.] Grado in cui un prodotto o sistema garantisce che i dati siano accessibili solo a chi è
    autorizzato ad accedervi.
  \item[Integrity.] Grado in cui un sistema, prodotto o componente impedisce modifiche non autorizzate di programmi o
    dati.
  \item[Authenticity.] Grado in cui si può dimostrare che l'identità di un soggetto o di una risorsa è quella
    dichiarata.
\end{description}
\end{definizione}

\figrif{fig:secm-tassonomia} riporta l'intera tassonomia e \tabrif{tab:secm-tassonomia} la riassume. Le percentuali
accanto a cinque smell indicano quanto ciascuno è citato nella rassegna.

\begin{figure}[ht]
\centering
\begin{tikzpicture}[
  s/.style={draw=rosso, text=rosso, minimum width=5.0cm, minimum height=0.5cm, align=center, font=\scriptsize, inner sep=2pt},
  r/.style={draw=verde, text=verde, minimum width=4.6cm, minimum height=0.5cm, align=center, font=\scriptsize, inner sep=2pt},
  e/.style={-{Stealth[length=1.6mm]}, thin, grigio}]
\node[font=\small\bfseries\sffamily, blu] at (-1.55,0.9) {Property};
\node[font=\small\bfseries\sffamily, rosso] at (3.4,0.9) {Smell};
\node[font=\small\bfseries\sffamily, verde] at (10.0,0.9) {Refactoring};
\foreach \i/\t in {1/Insufficient access control, 2/Publicly accessible microservices,
  3/Unnecessary privileges to microservices, 4/Own crypto code, 5/Non-encrypted data exposure,
  6/Hardcoded secrets, 7/Non-secured service-to-service communications, 8/Unauthenticated traffic,
  9/Multiple user authentication, 10/Centralized authorization}
  \node[s] (s\i) at (3.4,-0.75*\i+0.75) {\t};
\foreach \i/\t in {1/Use OAuth 2.0, 2/Add an API gateway, 3/Follow the least privilege principle,
  4/Use established encryption technologies, 5/Encrypt all sensitive data at rest, 6/Encrypt secrets at rest,
  7/Use mutual TLS, 8/Use OpenID Connect, 9/Use single sign-on, 10/Use decentralized authorization}
  \node[r] (r\i) at (10.0,-0.75*\i+0.75) {\t};
\foreach \i in {1,...,7,10} \draw[e] (s\i.east) -- (r\i.west);
\draw[e] (s8.east) -- (r7.west);
\draw[e] (s8.east) -- (r8.west);
\draw[e] (s9.east) -- (r8.west);
\draw[e] (s9.east) -- (r9.west);
\draw[e] (s9.east) -- (r2.west);
% proprieta
\draw[blu, fill=blu!10] (-2.35,0.25) rectangle (-1.95,-4.75);
\draw[blu, fill=blu!10] (-1.75,-1.25) rectangle (-1.35,-4.75);
\draw[blu, fill=blu!10] (-1.15,-2.0) rectangle (-0.75,-7.0);
\node[rotate=90, font=\scriptsize, text=blu] at (-2.15,-2.25) {Confidentiality};
\node[rotate=90, font=\scriptsize, text=blu] at (-1.55,-3.0) {Integrity};
\node[rotate=90, font=\scriptsize, text=blu] at (-0.95,-4.5) {Authenticity};
% percentuali
\foreach \i/\p in {1/43, 2/43, 7/55, 9/34, 10/33}
  \node[font=\scriptsize\bfseries, text=rosso, anchor=east] at ([xshift=-2pt]s\i.west) {\p\%};
\end{tikzpicture}
\caption{Tassonomia degli smell di sicurezza. Le barre a sinistra mostrano quali smell colpiscono ciascuna proprietà:
confidentiality i primi sette, integrity dal terzo al settimo, authenticity dal quarto al decimo. Ogni smell punta ai
refactoring che lo mitigano. Use mutual TLS e Use OpenID Connect servono due smell ciascuno; Add an API gateway serve
anche Multiple user authentication.}
\label{fig:secm-tassonomia}
\end{figure}

\begin{table}[ht]
\centering
\small
\caption{Smell di sicurezza, proprietà colpite e refactoring. C = confidentiality, I = integrity, A = authenticity.}
\label{tab:secm-tassonomia}
\begin{tabularx}{\linewidth}{>{\raggedright\arraybackslash}p{4.2cm}cc>{\raggedright\arraybackslash}X}
\toprule
Smell & Proprietà & \% & Refactoring \\
\midrule
Insufficient access control & C & 43 & use OAuth 2.0 \\
Publicly accessible microservices & C & 43 & add an API gateway \\
Unnecessary privileges to microservices & C, I & & follow the least privilege principle \\
Own crypto code & C, I, A & & use established encryption technologies \\
Non-encrypted data exposure & C, I, A & & encrypt all sensitive data at rest \\
Hardcoded secrets & C, I, A & & encrypt secrets at rest \\
Non-secured service-to-service communications & C, I, A & 55 & use mutual TLS \\
Unauthenticated traffic & A & & use mutual TLS, use OpenID Connect \\
Multiple user authentication & A & 34 & use single sign-on, add an API gateway, use OpenID Connect \\
Centralized authorization & A & 33 & use decentralized authorization \\
\bottomrule
\end{tabularx}
\end{table}

\subsection{Insufficient access control}
\label{sec:secm-accesso}

\textbf{Smell.} L'applicazione non applica il controllo degli accessi in alcuni microservizi.
\begin{itemize}
  \item Si apre la possibilità di un \termine{confused deputy problem}: un componente con privilegi legittimi viene
    indotto a usarli per conto di chi non li ha. Un attaccante può così ottenere dati che non dovrebbe vedere.
  \item Si rischia di violare la \textbf{confidentiality} dei dati e delle funzioni di business.
\end{itemize}
I permessi del client vanno verificati \emph{al momento della richiesta}. L'identità del client va verificata senza
introdurre latenza in più e contesa dovuta a chiamate frequenti a un servizio centralizzato.

\textbf{Refactoring}: \emph{use OAuth 2.0}, il più suggerito. OAuth 2.0 è un framework di sicurezza basato su token per
il controllo degli accessi delegato.
\begin{itemize}
  \item Il proprietario di una risorsa può concedere a un client l'accesso alla risorsa per suo conto.
  \item L'accesso è concesso per un tempo limitato e con un ambito (\ingl{scope}) limitato.
\end{itemize}

\subsection{Publicly accessible microservices}

\textbf{Smell.} Alcuni microservizi sono accessibili direttamente dai client esterni.
\begin{itemize}
  \item Ognuno di questi microservizi deve controllare autenticazione e autorizzazione per ogni richiesta.
  \item Le credenziali sono più esposte, e quindi aumenta la probabilità di violare la \textbf{confidentiality}.
  \item La manutenzione dell'applicazione costa di più.
\end{itemize}

\textbf{Refactoring}: \emph{add an API gateway} (\rif{sec:micro-motivi}), il più citato.
\begin{itemize}
  \item I microservizi non sono più accessibili direttamente dai client esterni.
  \item L'API gateway può imporre autenticazione, autorizzazione, \emph{throttling} e validazione del contenuto dei
    messaggi.
  \item I microservizi possono restare dietro un firewall.
\end{itemize}
Lo stesso refactoring risolve lo smell \emph{no API gateway} dei principi di design (capitolo~\ref{cap:smell}).

\subsection{Unnecessary privileges to microservices}

\textbf{Smell.} Ai microservizi sono concessi livelli di accesso, permessi o funzionalità che non servono a svolgere le
loro funzioni di business. Per esempio un servizio può leggere e scrivere dati in un database, o messaggi in una coda,
anche se non gli servono. Le risorse sono esposte senza necessità e la superficie d'attacco contro
\textbf{confidentiality} e \textbf{integrity} si allarga.

\textbf{Refactoring}: \emph{follow the least privilege principle}, il più citato.

\begin{definizione}[Principio del privilegio minimo]
Il \termine{least privilege principle} concede al codice in esecuzione solo i permessi necessari per completare i
compiti richiesti, e nient'altro.
\end{definizione}

\subsection{Own crypto code}

\textbf{Smell.} Codice crittografico ``fatto in casa''. Può causare problemi di \textbf{confidentiality},
\textbf{integrity} e \textbf{authenticity}. Può essere perfino peggio di nessuna cifratura, perché dà un falso senso di
sicurezza.

\textbf{Refactoring}: \emph{use established encryption technologies}.
\begin{itemize}
  \item Usare librerie di cifratura collaudate a fondo dalla comunità, riviste e corrette regolarmente.
  \item Evitare algoritmi di cifratura sperimentali, che possono avere vulnerabilità non ancora note.
\end{itemize}

\subsection{Non-encrypted data exposure}

\textbf{Smell.} L'applicazione espone per sbaglio dati sensibili. Per esempio i dati sono memorizzati senza cifratura,
oppure la protezione usata ha vulnerabilità o difetti. Gli intrusi possono leggere e modificare i dati, comprese le
credenziali: \textbf{confidentiality}, \textbf{integrity} e \textbf{authenticity} possono essere violate.

\textbf{Refactoring}: \emph{encrypt all sensitive data at rest}. Tutti i dati sensibili devono restare sempre cifrati
ed essere decifrati solo quando servono (\rif{sec:sic-dati}).
\begin{itemize}
  \item La maggior parte dei DBMS supporta la cifratura automatica.
  \item I sistemi operativi supportano la cifratura a livello di disco.
  \item La cifratura a livello di applicazione è un'altra possibilità.
  \item Vanno cifrati anche i dati nella cache.
\end{itemize}
La cifratura consuma risorse: si individuano i dati critici e si cifra solo ``quanto serve''.

\subsection{Hardcoded secrets}

\textbf{Smell.} Segreti (chiavi delle API, segreti dei client, credenziali) scritti direttamente:
\begin{itemize}
  \item nel codice sorgente del microservizio;
  \item negli script di distribuzione dell'applicazione, per esempio variabili d'ambiente che passano segreti in una
    specifica Docker Compose (\rif{sec:container-compose}).
\end{itemize}
I segreti non vanno mai memorizzati nelle variabili d'ambiente. Possono essere esposti per sbaglio: per esempio un
gestore delle eccezioni può mandare queste informazioni a una piattaforma di logging. Così la \textbf{confidentiality} e
l'\textbf{integrity} dei segreti possono essere violate.

\textbf{Refactoring}: \emph{encrypt secrets at rest}, cioè cifrare i segreti. Le buone pratiche sono due.
\begin{itemize}
  \item Non conservare le credenziali insieme alle applicazioni o nei repository del codice sorgente.
  \item Non usare variabili d'ambiente per passare i segreti.
\end{itemize}

\subsection{Non-secured service-to-service communications}

\textbf{Smell.} I microservizi interagiscono senza usare un canale di comunicazione sicuro. I dati sono esposti ad
attacchi \emph{man-in-the-middle} (un intruso si inserisce tra i due servizi), a intercettazioni (\ingl{eavesdropping})
e a manomissioni (\ingl{tampering}). Gli intrusi possono leggere e alterare i dati in transito: \textbf{confidentiality},
\textbf{integrity} e \textbf{authenticity} sono a rischio.

\textbf{Refactoring}: \emph{use mutual TLS}. Nel \termine{mutual TLS} entrambe le parti, non solo il server, presentano
un certificato (\rif{sec:sic-tls}).
\begin{itemize}
  \item È una soluzione largamente accettata: cifra i dati in transito e ne garantisce integrità e riservatezza.
  \item Permette anche a un microservizio di identificare con certezza il microservizio con cui parla
    (\emph{autenticazione reciproca}).
\end{itemize}

\subsection{Unauthenticated traffic}

\textbf{Smell.} Richieste non autenticate, dall'esterno o tra microservizi. È essenziale che i microservizi possano
autenticarsi a vicenda, soprattutto quando si passa il contesto utente. Se il traffico non è autenticato, i
microservizi sono esposti ad attacchi. Esempi sono la manomissione dei dati, il denial of service (\rif{sec:sic-dos}) e
l'\emph{elevation of privileges}, cioè l'acquisizione di privilegi non dovuti. La proprietà a rischio è
l'\textbf{authenticity}.

\textbf{Refactoring}: due.
\begin{itemize}
  \item \emph{Use mutual TLS}.
  \item \emph{Use OpenID Connect}. \termine{OpenID Connect} usa JWT che contengono le informazioni dell'utente
    autenticato. I microservizi verificano l'identità dell'utente presso un \emph{authorization server}.
\end{itemize}

\subsection{Multiple user authentication}

\textbf{Smell.} L'applicazione offre più punti di accesso per l'autenticazione degli utenti.
\begin{itemize}
  \item Ogni punto di accesso è un possibile vettore d'attacco, con cui un intruso può autenticarsi come utente finale:
    l'\textbf{authenticity} è a rischio.
  \item La manutenzione dell'applicazione costa di più.
\end{itemize}

\textbf{Refactoring}: \emph{use single sign-on}. Con il \termine{single sign-on} c'è un unico punto d'ingresso che
gestisce l'autenticazione degli utenti e impone la sicurezza a tutte le loro richieste. Rende anche più semplici la
conservazione dei log e i controlli (\ingl{auditing}). Il single sign-on si ottiene:
\begin{itemize}
  \item aggiungendo un API gateway, se non c'è già;
  \item usando OpenID Connect, per condividere il contesto utente.
\end{itemize}

\subsection{Centralized authorization}

L'autorizzazione si può imporre ai margini dell'applicazione (nell'API gateway), in ogni microservizio, o in entrambi.

\textbf{Smell.} L'applicazione gestisce l'autorizzazione solo ai margini.
\begin{itemize}
  \item Il punto ``centrale'' di autorizzazione diventa un collo di bottiglia: prestazioni ed efficienza peggiorano.
  \item C'è la possibilità di un'altra forma di confused deputy problem (\rif{sec:secm-accesso}): i microservizi si
    fidano del gateway solo in base alla sua identità. L'\textbf{authenticity} può essere violata.
\end{itemize}

\textbf{Refactoring}: \emph{use decentralized authorization}.
\begin{itemize}
  \item Si trasmette un token di accesso (per esempio un JWT) insieme a ogni richiesta a un microservizio.
  \item Il microservizio concede l'accesso al chiamante solo se riceve un token noto e corretto.
\end{itemize}

Riferimento: F. Ponce, J. Soldani, H. Astudillo, A. Brogi, \emph{Smells and Refactorings for Microservices Security:
A Multivocal Literature Review}, The Journal of Systems and Software, 2022.

\section{Come eliminare gli smell di sicurezza}
\label{sec:secm-eliminare}

\begin{algoritmo}[Eliminare gli smell di sicurezza]
\begin{enumerate}
  \item \textbf{Detect}: individuare gli smell di sicurezza.
  \item \textbf{Triage}: decidere quali smell affrontare e in che ordine, tenendo conto di:
    \begin{itemize}
      \item quanto è rilevante il servizio;
      \item l'impatto dello smell sugli attributi di qualità;
      \item lo sforzo necessario per il refactoring;
      \item l'eventuale necessità di interagire con altri team.
    \end{itemize}
  \item \textbf{Refactor}: scegliere e applicare un refactoring.
\end{enumerate}
\end{algoritmo}

Il triage ha più dimensioni da bilanciare insieme. Per questo servono \textbf{strumenti di supporto}, come quelli per
gli smell architetturali (\rif{sec:smell-toolchain}).

\begin{puntichiave}
\begin{itemize}
  \item Sette sfide:
    \begin{itemize}
      \item superficie d'attacco più ampia;
      \item controlli distribuiti costosi (trust-the-network contro zero trust);
      \item fiducia tra servizi da automatizzare (certificati);
      \item tracing difficile;
      \item container immutabili e credenziali da ruotare;
      \item contesto utente da passare (JWT);
      \item responsabilità divise tra team.
    \end{itemize}
  \item Una multivocal review di 58 riferimenti individua 10 smell, 10 refactoring e 3 proprietà (confidentiality,
    integrity, authenticity, dallo standard ISO/IEC 25010).
  \item Controllo degli accessi: insufficient access control $\to$ OAuth 2.0; publicly accessible microservices $\to$
    API gateway; unnecessary privileges $\to$ least privilege principle.
  \item Cifratura: own crypto code $\to$ tecnologie consolidate; dati e segreti non cifrati $\to$ cifrarli a riposo;
    comunicazioni non protette $\to$ mutual TLS.
  \item Autenticazione e autorizzazione: unauthenticated traffic $\to$ mutual TLS, OpenID Connect; multiple user
    authentication $\to$ single sign-on (API gateway + OpenID Connect); centralized authorization $\to$ autorizzazione
    decentralizzata con token.
  \item Per eliminare gli smell: detect, triage, refactor, con strumenti di supporto.
\end{itemize}
\end{puntichiave}

\begin{domande}
\begin{enumerate}
  \item Perché i microservizi allargano la superficie d'attacco rispetto a un monolite?
  \item Che differenza c'è tra trust-the-network e zero-trust networking?
  \item Perché stabilire la fiducia tra microservizi richiede automazione?
  \item Perché tracciare le richieste è più difficile con i microservizi? Che differenza c'è tra log e trace?
  \item Perché i container complicano la gestione di credenziali e politiche?
  \item Come si condivide il contesto utente tra microservizi?
  \item Definire confidentiality, integrity e authenticity secondo ISO/IEC 25010.
  \item Che cos'è il confused deputy problem? In quali smell compare?
  \item Perché i segreti non vanno passati con le variabili d'ambiente?
  \item Che cosa aggiunge il mutual TLS al TLS? Quali smell risolve?
  \item Come si realizza il single sign-on?
  \item Quali passi servono per eliminare gli smell di sicurezza? Su che cosa si basa il triage?
\end{enumerate}
\end{domande}
